\chapter{The First Variation}
% OUTLINE Ch1 "The First Variation (geometry's face: the extremal
% structure)". Laws: INTUITION LAW (every picture in the \intuition env,
% adds-no-claim), CITATION (Bernoulli 1696 challenge + the Newton night
% at Cavendish grade; Fermat 1662; Euler 1744; Lagrange 1755/1788 w/ the
% Euler-adoption attribution story; Vol 2 Ch7 geodesics re-felt), the
% teleology FENCE (best = property of the functional, not purpose in the
% world). NO device appearance (lens sites are Ch2+Ch4). Gate: Kodo.

\begin{intuition}
Dip a bent wire in soap solution and lift it out: the film that spans it
is the least area the wire permits. Hang a chain from two nails: the
curve it settles into is the one that puts its weight lowest. Watch light
cross from air into water: of all the bent paths it might take, it takes
the one that spends the least time. Three benches, one feeling --- that
nature, offered every candidate, somehow picks the best one --- and the
feeling is worth having before it is examined, because a whole face of
mathematics is the examination.
\end{intuition}

The examination began as a public dare. In June of 1696, Johann Bernoulli
posed a challenge to ``the sharpest mathematicians in the whole world'':
given two points, one above the other but not vertically so, find the
curve down which a bead slides under gravity in the least time. The
answer is not the straight line --- the straight line is shortest but not
fastest --- and the problem, the \emph{brachistochrone}, is the subject's
founding legend for what happened next. Newton, by the account that has
come down through his niece, returned home from the Mint at four in the
afternoon, found the challenge, and had solved it by four the next
morning; he published the solution anonymously, and Bernoulli, reading
it, is said to have remarked that he knew the lion by his claw. The
curve is the cycloid, the same curve traced by a point on a rolling
wheel, and five solutions arrived in all --- both Bernoullis, Leibniz,
l'H\^opital, and the anonymous lion. What matters to this chapter is not
the roster but the kind of question: not \emph{what is the value of this
quantity} but \emph{which curve, among all candidates, is extreme} ---
and a question quantified over all candidates needs a new kind of proof.

The proof structure that answers it was built in two generations. Euler,
in his treatise of 1744, gave the first systematic method and the
equation that bears the subject's fingerprint: a candidate curve is
extremal exactly when a certain differential condition holds along it.
Then Lagrange --- nineteen years old, writing to Euler in 1755 ---
replaced Euler's geometric derivation with the analytic device the
subject still uses, the variation itself: perturb the candidate by an
arbitrary small admissible wiggle, expand to first order, and demand
that the first-order change vanish for \emph{every} admissible wiggle.
The attribution deserves its sentence, in this series' tradition: Euler,
receiving the young man's method, abandoned his own, adopted Lagrange's,
gave the subject the name it still carries --- the calculus of
variations --- and said whose the method was. Lagrange's mature
statement is the \emph{M\'ecanique analytique} of 1788; Fermat's least
time, which the soap-and-light intuition above leans on, was stated for
optics in 1662 and stands as the structure's great anticipation.

Now say the structure plainly, because it is the chapter's cargo and the
face's signature. To prove a candidate extremal, one does not compare it
with each rival --- the rivals are uncountable. One \emph{varies} it:
adds an arbitrary admissible perturbation, differentiates the quantity
of interest along that perturbation, and shows the derivative vanishes
at the candidate, whatever the perturbation was. The universal
quantifier --- \emph{for every} wiggle --- is carried not by exhaustion
but by the arbitrariness of one symbol, and the vanishing condition
condenses, by the Euler--Lagrange argument, into a differential
equation checked pointwise along the candidate. That is the first
variational proof structure: a claim about a whole ocean of candidates,
proved by one derivative taken in a direction left arbitrary. The reader
has owned its most important instance since the last volume: the
geodesics of Chapter 7 there --- the locally-shortest paths the ruler
selects --- are exactly such theorems. There, they were objects; here,
the reader is asked to re-feel them as \emph{proofs} --- a geodesic is a
statement whose demonstration is a vanishing first variation --- and
that re-feeling is what this volume means by its title.

\begin{intuition}
The feeling that carries the structure: stand on the candidate and
jiggle. If every jiggle, to first order, costs nothing and gains
nothing, you are standing somewhere special. The proof is the jiggling
made honest --- one arbitrary jiggle standing for all of them.
\end{intuition}

And now the fence, because the opening feeling --- \emph{nature picks
the best} --- contains a word this book must not let harden. ``Best'' is
a property of the functional some person chose to extremize: least area
for the film, least potential energy for the chain, least time for the
ray. The mathematics proves that the observed configuration extremizes
that chosen quantity; it does not prove that nature \emph{prefers},
\emph{intends}, or \emph{deliberates}. The light ray does not survey its
options; the variational statement and the pointwise differential law
are two spellings of one content --- Lagrange's own treatise exists to
translate between them --- and the appearance of choice is the reader's,
imported with the word ``best.'' The intuition remains a fine way to
find the functional and to feel the proof; teleology is what it must
not become, and this paragraph is the label saying so.

What the first variation cannot see is the chapter's honest limit, and
the next chapter's opening. A vanishing first derivative marks the flat
places --- but the mountain pass is as flat, at its saddle, as the
valley floor. Which flat places are stable, which are passes, and what
the answer has to do with the \emph{shape} of the space the candidates
live in --- that question needs the second variation, and it carries
this volume to its second face.
